
\section{VANET Architecture and Communication Technologies}
\label{sec:2}

Vehicular Ad Hoc Networks (VANETs) constitute a
specialized class of wireless mobile networks
designed to support communication among vehicles,
roadside infrastructure, and other road users.
Their architecture combines mobile communication
nodes, fixed infrastructure, wireless access
technologies, and application-level services
\cite{Karagiannis2011}.

Unlike conventional mobile ad hoc networks,
vehicular networks operate in environments
constrained by road geometry, traffic regulations,
vehicle mobility, and rapidly changing wireless
connectivity.

The design of a VANET must accommodate diverse
communication requirements. Safety applications
typically prioritize timely and reliable
information dissemination, whereas traffic
management and information services may require
broader coverage, higher data capacity, or
infrastructure connectivity. Consequently,
vehicular communication architecture must
consider mobility, scalability, interoperability,
and security across multiple system layers
\cite{Karagiannis2011}.

\subsection{Fundamental VANET Architecture}
\label{sec:2A}

A typical VANET consists of vehicles equipped
with communication capabilities, roadside
communication equipment, and supporting
transportation or network infrastructure.
Vehicles can exchange information directly
or through infrastructure-assisted connections
\cite{Hartenstein2008}.

\subsubsection{On-Board Units}

An On-Board Unit (OBU) is a communication
component installed inside a vehicle.
It provides wireless connectivity and
supports the transmission and reception
of vehicular messages. Depending on the
system design, an OBU may interact with
vehicle sensors, positioning systems,
processing units, and in-vehicle
communication interfaces.

OBUs can support periodic awareness
messages as well as event-triggered
notifications. Periodic messages provide
information about vehicle state, such
as position, speed, and movement
direction, while event-triggered
messages may communicate hazardous
conditions or unexpected road events.

The effectiveness of an OBU depends
on several factors, including wireless
channel conditions, transmission
parameters, positioning accuracy,
message processing capability,
and security mechanisms
\cite{Zeadally2012}.

\subsubsection{Roadside Units}

Roadside Units (RSUs) are fixed
communication nodes deployed at
selected locations along transportation
infrastructure, including intersections,
highways, and urban roads.

An RSU may exchange information
with nearby vehicles, distribute
traffic-related notifications,
and provide connectivity to
transportation management systems
or external communication networks.

RSU deployment introduces important
design considerations. Increasing
the number of roadside units may
improve infrastructure availability,
but it also increases deployment
and maintenance costs. The appropriate
deployment density depends on
road characteristics, traffic
conditions, and service requirements
\cite{Karagiannis2011}.

\subsubsection{Application Units and Vehicle Sensors}

Application Units (AUs) represent
the processing or software components
responsible for providing vehicular
services. Depending on the architecture,
these functions may be integrated
with the OBU or implemented using
separate in-vehicle systems.

Vehicle sensors and positioning
technologies can provide information
about movement, location, and
surrounding conditions. This
information may support applications
such as collision warnings,
traffic monitoring, and
cooperative awareness.

However, the accuracy and
availability of sensor information
can influence the quality of
vehicular communication services.
Incorrect or outdated information
may reduce the usefulness of
safety-related messages
\cite{Zeadally2012}.

\subsection{Communication Architecture}
\label{sec:2B}

Vehicular communication is commonly
described through several interaction
models. These models differ according
to the communicating entities,
infrastructure requirements,
and intended applications.

\subsubsection{Vehicle-to-Vehicle Communication}

Vehicle-to-Vehicle (V2V) communication
allows vehicles to exchange information
directly with nearby vehicles through
compatible wireless interfaces.

V2V communication is particularly
relevant to cooperative awareness
and road safety applications.
For example, a vehicle detecting
emergency braking may transmit
information to neighboring vehicles,
allowing them to become aware
of a potential hazard.

Direct communication can reduce
dependence on centralized
infrastructure for local
information exchange. However,
its performance is influenced
by communication range,
channel congestion, obstacles,
and the relative movement
of vehicles \cite{Hartenstein2008}.

\subsubsection{Vehicle-to-Infrastructure Communication}

Vehicle-to-Infrastructure (V2I)
communication enables information
exchange between vehicles and
roadside communication systems.

Through V2I communication,
vehicles may receive traffic
updates, road condition
information, or infrastructure
messages. Roadside systems
may also collect information
from connected vehicles
to support traffic management.

Unlike direct V2V communication,
V2I services depend on the
availability and capabilities
of the deployed infrastructure.
Coverage limitations and
deployment costs are therefore
important considerations
\cite{Karagiannis2011}.

\subsubsection{Vehicle-to-Pedestrian Communication}

Vehicle-to-Pedestrian (V2P)
communication involves information
exchange between vehicles and
pedestrians or other vulnerable
road users equipped with
compatible communication devices.

Potential applications include
warnings about approaching
vehicles and increased awareness
at intersections or pedestrian
crossings.

Practical V2P systems must
consider device energy consumption,
positioning uncertainty,
communication availability,
and privacy requirements.

\subsubsection{Vehicle-to-Network Communication}

Vehicle-to-Network (V2N)
communication provides connectivity
between vehicles and external
network services, commonly through
cellular infrastructure.

V2N connectivity can support
cloud-based traffic information,
remote services, and access
to transportation management
platforms.

Unlike direct local communication,
network-assisted services may
depend on cellular coverage,
network congestion, and
end-to-end service availability
\cite{Karagiannis2011}.

\begin{table}[htbp]
\centering
\caption{Comparison of VANET Communication Models}
\label{tab:communication_models}
\footnotesize
\begin{tabular}{|p{1.0cm}|p{2.0cm}|p{3.4cm}|}
\hline
\textbf{Mode} &
\textbf{Entities} &
\textbf{Typical Applications} \\
\hline
V2V &
Vehicle--Vehicle &
Cooperative awareness and hazard warnings \\
\hline
V2I &
Vehicle--Infrastructure &
Traffic information and roadside services \\
\hline
V2P &
Vehicle--Pedestrian &
Vulnerable road-user awareness \\
\hline
V2N &
Vehicle--Network &
Cloud services and network-assisted information \\
\hline
\end{tabular}
\end{table}

\subsection{Wireless Communication Technologies}
\label{sec:2C}

Different wireless technologies have
been developed or adapted to support
vehicular communication. Their
characteristics vary in terms
of channel access, network
architecture, coverage, and
communication requirements
\cite{Karagiannis2011}.

\subsubsection{IEEE 802.11p and DSRC}

IEEE 802.11p introduced amendments
to wireless local area networking
for operation in vehicular
environments. It has been widely
associated with Dedicated
Short-Range Communications (DSRC)
and research on direct vehicular
communication.

The technology supports wireless
information exchange between
compatible devices without
requiring a conventional
access-point association process
for every local interaction.

However, communication performance
may be affected by contention
for shared wireless resources,
especially when many vehicles
attempt to transmit simultaneously.

The IEEE 802.11p-related approach
has also influenced vehicular
communication systems based
on the European ITS-G5 framework
\cite{Karagiannis2011}.

\subsubsection{Cellular Vehicle-to-Everything}

Cellular Vehicle-to-Everything
(C-V2X) provides an alternative
family of technologies for
vehicular communication.

LTE-V2X includes a direct
communication interface commonly
known as PC5, as well as
cellular network connectivity
through the conventional
network interface.

Direct communication can support
local information exchange,
while network-assisted
communication enables access
to broader services.

C-V2X systems must address
resource allocation, interference,
mobility, and coexistence
with other wireless services.

\subsubsection{5G NR-V2X}

New Radio Vehicle-to-Everything
(NR-V2X) extends cellular
vehicular communication
capabilities using technologies
associated with 5G New Radio.

NR-V2X is designed to support
a broader range of vehicular
use cases, including advanced
cooperative driving scenarios.

The requirements of these
applications may involve
high reliability, low latency,
and the exchange of more
complex information between
connected road users.

Nevertheless, performance
depends on deployment
conditions, network
configuration, and
application requirements.

\begin{table}[htbp]
\centering
\caption{Qualitative Comparison of Vehicular Wireless Technologies}
\label{tab:wireless_technologies}
\footnotesize
\begin{tabular}{|p{1.7cm}|p{2.2cm}|p{2.5cm}|}
\hline
\textbf{Technology} &
\textbf{Main Approach} &
\textbf{Key Consideration} \\
\hline
IEEE 802.11p &
Direct wireless access &
Shared-channel contention \\
\hline
LTE-V2X &
Cellular-based direct and network communication &
Radio resource management \\
\hline
NR-V2X &
Advanced cellular vehicular communication &
Support for demanding V2X use cases \\
\hline
\end{tabular}
\end{table}

\subsection{Communication Protocol Architecture}
\label{sec:2D}

Vehicular communication systems
use protocol architectures
that organize functions across
different layers. Although
specific implementations vary,
communication generally involves
physical transmission,
medium access, networking,
transport, facilities,
and application-level services.

\subsubsection{Physical Layer}

The physical layer is responsible
for transmitting information
over wireless channels.
Its operation is affected
by signal propagation,
interference, mobility,
and environmental obstacles
\cite{Zeadally2012}.

Vehicular wireless channels
may experience rapid variations
because the transmitter,
receiver, or surrounding
objects are moving.

\subsubsection{Medium Access Control}

The Medium Access Control
(MAC) layer manages access
to shared wireless resources.

In vehicular environments,
multiple nodes may attempt
to transmit messages over
the same channel. Effective
channel access mechanisms
are therefore important
for reducing collisions,
contention, and delays
\cite{Zeadally2012}.

\subsubsection{Networking and Transport}

Networking mechanisms support
message forwarding and
information dissemination
between communication nodes.

Depending on the application,
information may be distributed
through direct broadcast,
geographic forwarding,
or infrastructure-assisted
communication.

Transport mechanisms and
application protocols must
also account for differences
between time-critical
safety messages and
less urgent information
services \cite{Karagiannis2011}.

\subsubsection{Facilities and Applications}

The facilities and application
layers support the generation,
processing, and interpretation
of vehicular information.

Examples include cooperative
awareness information and
event-driven hazard
notifications. Such services
depend on suitable message
formats, processing mechanisms,
and security protections.

\subsection{Mobility and Network Characteristics}
\label{sec:2E}

The mobility characteristics
of vehicles distinguish VANETs
from many other wireless networks.
Vehicle movement is constrained
by road geometry, traffic signals,
speed limits, and driving behavior.

\textbf{Dynamic Topology:}
As vehicles enter and leave
communication range, network
connections change over time.
This creates difficulties
for maintaining stable routes.

\textbf{Variable Vehicle Density:}
Sparse traffic may result
in disconnected network
segments, whereas dense
traffic may increase
wireless channel contention.

\textbf{Urban and Highway Environments:}
Urban roads contain intersections,
buildings, and obstacles
that can affect propagation.
Highway scenarios often
involve higher speeds and
rapid changes in relative
vehicle positions.

\textbf{Short Link Duration:}
The duration of a direct
wireless link depends on
relative mobility,
communication range,
and propagation conditions.
Short-lived links may
reduce the effectiveness
of routing mechanisms
that require stable
end-to-end paths.

These characteristics
demonstrate why communication
protocols designed for
conventional networks
may require adaptation
before being applied
to vehicular environments
\cite{Hartenstein2008,Zeadally2012}.

\subsection{Architecture Design Considerations}
\label{sec:2F}

A reliable VANET architecture
must balance several
potentially competing
design objectives.

First, safety-related
communication requires
timely message delivery
and dependable information
exchange. Second, large
vehicular populations
require scalable
communication mechanisms
that limit unnecessary
channel usage.

Third, interoperability
between vehicles,
infrastructure, and
communication technologies
is important for
practical deployment.
Finally, authentication,
message integrity,
and privacy protection
must be considered
throughout the system.

The selection of
communication technology
and network architecture
should therefore depend
on the intended application,
mobility environment,
deployment constraints,
and required level
of communication reliability
\cite{Karagiannis2011,Zeadally2012}.
